\begin{textsample}{POS Dim 3 – mt – Score 29.00 – mt/mt\_inf\_21.txt}  \label{ex:f3_pos_074}
2.4 \textbf{Escuta}, Fala, Pensamento e Imaginação ( EF ) Esse campo envolve a oralidade, a \textbf{escuta}, o pensamento e a imaginação, que devem ser estimulados desde a Educação \textbf{Infantil}.
Na \textbf{pequena} \textbf{infância}, a aquisição e o domínio da linguagem verbal está vinculada à constituição do pensamento, à fruição literária, sendo também instrumento de apropriação dos demais conhecimentos.
Nesse sentido, a prática pedagógica nas \textbf{instituições} educativas deve prever espaços, tempos, materiais e experiências que privilegiam as interações, para que as \textbf{crianças} possam se expressar, imaginar, criar, comunicar, organizar pensamentos e ideias, propiciando -lhes o acesso aos conhecimentos científicos e atividades produzidas coletivas e historicamente por diferentes grupos sociais, bem como a ampliação das suas experiências culturais.
É função da primeira etapa da Educação Básica, mais do que inserir, dar acesso ou entrada na cultura escrita, propiciar às \textbf{crianças} a aproximação, a participação na cultura escrita ou na “ cultura dos escritos ” ( GALVÃO, 2016 ).
Entende -se por cultura escrita, como o lugar que o escrito ( todo e qualquer evento ou prática que tenha como mediação a palavra escrita ) ocupa em determinada sociedade, comunidade ou grupo social.
Cabe à escola considerar que a Educação \textbf{Infantil} não tem compromisso com uma proposta de alfabetização.
Muito mais importante que ensinar as letras do alfabeto é familiarizar as \textbf{crianças}, desde \textbf{bebês}, com práticas sociais em que a leitura e a escrita estejam presentes exercendo funções diversas nas interações sociais.
A BNCC assim apresenta esse campo de experiência : > Desde o \textbf{nascimento}, as \textbf{crianças} participam de situações comunicativas cotidianas com as pessoas com as quais interagem.
As primeiras formas de interação do \textbf{bebê} são os movimentos do seu corpo, o olhar, a postura corporal, o sorriso, o choro e outros recursos vocais, que ganham sentido com a interpretação do outro.
Progressivamente, as \textbf{crianças} vão ampliando e enriquecendo seu vocabulário e demais recursos de expressão e de compreensão, apropriando -se da língua materna – que se torna, pouco a pouco, seu veículo privilegiado de interação.
Na Educação \textbf{Infantil}, é importante promover experiências nas quais as \textbf{crianças} possam falar e \textbf{ouvir}, potencializando sua participação na cultura oral, pois é na \textbf{escuta} de histórias, na participação em conversas, nas descrições, nas narrativas elaboradas individualmente ou em grupo e nas implicações com as múltiplas linguagens que a \textbf{criança} se constitui ativamente como sujeito singular e pertencente a um grupo social.
Desde cedo, a \textbf{criança} manifesta \textbf{curiosidade} com relação à cultura escrita : ao \textbf{ouvir} e acompanhar a leitura de textos, ao observar os muitos textos que circulam no contexto familiar, comunitário e escolar, ela vai construindo sua concepção de língua escrita, reconhecendo diferentes usos sociais da escrita, dos gêneros, suportes e portadores.
Na Educação \textbf{Infantil}, a \textbf{imersão} na cultura escrita deve partir do que as \textbf{crianças} conhecem e das \textbf{curiosidades} que deixam transparecer.
As experiências com a literatura \textbf{infantil}, propostas pelo educador, mediador entre os textos e as \textbf{crianças}, contribuem para o desenvolvimento do gosto pela leitura, do \textbf{estímulo} à imaginação e da ampliação do conhecimento de mundo.
Além disso, o contato com histórias, contos, fábulas, poemas, cordéis etc. propicia a familiaridade com \textbf{livros}, com diferentes gêneros literários, a diferenciação entre ilustrações e escrita, a aprendizagem da direção da escrita e as formas corretas de manipulação de \textbf{livros}.
Nesse convívio com textos escritos, as \textbf{crianças} vão construindo hipóteses sobre a escrita que se revelam, inicialmente, em rabiscos e garatujas e, à medida que vão conhecendo letras, em escritas espontâneas, não convencionais, mas já \textbf{indicativas} da compreensão da escrita como sistema de representação da língua ( BRASIL, 2017, p. 40 ). 2.4.1 Direitos de aprendizagem e desenvolvimento no campo de experiências \textbf{Escuta}, Fala, Pensamento e Imaginação No campo de experiências \textbf{escuta}, fala, pensamento e imaginação os direitos de aprendizagem podem ser assegurados quando oportunizamos as \textbf{crianças} : - CONVIVER com \textbf{crianças} e \textbf{adultos} em situações comunicativas cotidianas, estabelecendo modos de pensar, falar, imaginar, sentir, narrar e conhecer. - \textbf{BRINCAR} com diferentes gêneros textuais tanto na ordem do narrar, argumentar, expor, relatar e descrever ( parlendas, trava-línguas, adivinhas, memória, \textbf{rodas} e \textbf{brincadeiras} cantadas, jogos e imagens, entre outros ), expandido o repertório das manifestações culturais locais e de outras culturas, enriquecendo sua linguagem oral, corporal, musical, \textbf{dramática}, escrita, dentre outras. - PARTICIPAR de momentos que envolvem as \textbf{rodas} de conversa, relatos de experiências, contação e leitura de histórias e poesias, produção de narrativas, representação de papéis no faz de conta, exploração de materiais impressos e de variedades linguísticas, construindo diversas formas de organizar o pensamento. - EXPLORAR diversos \textbf{sons}, \textbf{gestos}, expressões, imagens, textos escritos, além dos sentidos das palavras, nos diferentes gêneros literários apropriando -se desses elementos para elaborar novas falas, enredos, histórias e escritas convencionais ou não. - EXPRESSAR sentimentos, ideias, \textbf{percepções}, \textbf{desejos}, necessidades, pontos de vista, informações, dúvidas, descobertas, empregando múltiplas linguagens, considerando as manifestações e opiniões dos \textbf{colegas} e \textbf{adultos}. - CONHECER -SE e reconhecer suas preferências para as \textbf{brincadeiras}, lugares, histórias, autores, gêneros linguísticos, e seu interesse em produzir com a linguagem verbal.
FONTE : BRASIL ( 2018 ) 2.4.2 Objetivos de aprendizagem e desenvolvimento no Campo de Experiências \textbf{Escuta}, Fala, Pensamento e Imaginação A partir dessa caracterização do campo de experiência \textbf{escuta}, fala, pensamento e imaginação, foram estabelecidos objetivos de aprendizagem e desenvolvimento, conforme o quadro a seguir : 2.4.3 Professor, o que garantir no planejamento?
As práticas pedagógicas que compõem o Documento de Referência Curricular para a Educação \textbf{Infantil} devem garantir experiências que “ favoreçam a \textbf{imersão} das \textbf{crianças} nas diferentes linguagens e o progressivo domínio por elas de vários gêneros e formas de experiência : gestual, verbal, \textbf{plástica}, \textbf{dramática} e musical ”, conforme disposto nas DCNEI ( 2009 ), artigo 9º, inciso II ( BRASIL, 2009 ).
Para tanto, é necessário que o planejamento seja elaborado a partir de reflexões que os direitos de aprendizagem e desenvolvimento possibilitam e devem contemplar práticas que propiciem a \textbf{bebês}, \textbf{crianças} bem \textbf{pequenas} e \textbf{crianças} \textbf{pequenas} “ experiências de narrativas, de apreciação e interação com a linguagem oral e escrita, e convívio com diferentes suportes e gêneros textuais, orais e escritos ”, como disposto nas DCNEI ( 2009 ), artigo 9º, inciso II ( BRASIL, 2009 ).
É imprescindível que o professor \textbf{acolha} as necessidades, os \textbf{desejos} e as manifestações das \textbf{crianças}, suas histórias de vida, a realidade de suas famílias e o contexto no qual estão inseridas, e as assumam como produtoras de cultura, como aquelas que \textbf{inventam} o mundo, com uma história de cultura a serem partilhadas.
Isso implica em planejar o cotidiano levando em conta o ponto de vista das \textbf{crianças}, seu jeito de conhecer e interagir com o mundo à sua \textbf{volta}, seus modos de se expressar por meio de diferentes linguagens, movimentos e produções ” ( KRAMER e BARBOSA, 2016 ).
Assim, na perspectiva da invenção e da expressão por meio das diferentes linguagens, é indispensável garantir tempo : o tempo para falar, \textbf{ouvir}, \textbf{brincar}, ler histórias, desenhar, dentro e fora das salas, comer, descansar, \textbf{escutar} as \textbf{crianças} abrindo espaço para suas manifestações, e também promover o contato com o conhecimento científico e cultural, com a arte e as culturas.
Dessa forma, propiciar a \textbf{imersão} das \textbf{crianças} na cultura escrita.

% matched lemmas: acolher, adulto, bebê, brincadeira, brincar, colega, criança, curiosidade, desejo, dramático, escuta, escutar, estímulo, gesto, imersão, indicativo, infantil, infância, instituição, inventar, livro, nascimento, ouvir, pequeno, percepção, plástico, roda, som, volta
\end{textsample}
